\chapter{Hyperbolic manifolds} \label{varieta:chapter} \label{manifold:chapter} 
 \label{orbifold:chapter} \label{orbifolds:chapter}

A hyperbolic manifold is a Riemannian manifold locally isometric to the hyperbolic space $\matH^n$. Maybe the most striking aspect of geometric topology is that, despite this quite restrictive definition, there are plenty of hyperbolic manifolds around, especially in the dimensions $n=2$ and $3$. For that reason hyperbolic manifolds (and hence hyperbolic geometry) play a central role in the topology of surfaces and three-manifolds.

The study of complete hyperbolic manifolds is tightly connected to that of discrete subgroups in the Lie group $\Iso(\matH^n)$ and of polyhedra in $\matH^n$, so it has both an algebraic and geometric flavour. We start this chapter by describing these connections; then we show some examples and discuss some important variations: non-complete hyperbolic manifolds, hyperbolic manifolds with geodesic boundary, cone manifolds, and orbifolds.

\section{Discrete groups of isometries}
We define hyperbolic manifolds and prove a crucial theorem, that says that every complete hyperbolic manifold is isometric to a quotient $\matH^n/_\Gamma$ for some discrete group $\Gamma <\Iso(\matH^n)$ acting freely on $\matH^n$.

\subsection{Hyperbolic, flat, and elliptic manifolds}
We introduce three important classes of Riemannian manifolds.
\index{hyperbolic manifold} \index{flat manifold} \index{elliptic manifold}

\begin{defn} A \emph{hyperbolic} (resp.~\emph{flat} or \emph{elliptic}) \emph{manifold} is a connected Riemannian $n$-manifold that may be covered by open sets isometric to open sets of $\matH^n$ (resp.~$\matR^n$ o $S^n$). 
\end{defn}

A hyperbolic (resp.~flat or elliptic) manifold has constant sectional curvature $-1$ (resp.~$0$ or $+1$). We show that the model $\matH^n$ is indeed unique.


\begin{teo} \label{simply:teo}
Every complete simply connected hyperbolic $n$-manifold $M$ is isometric to $\matH^n$.
\end{teo}
\begin{proof}
Pick a point $x\in M$ and choose an isometry $D\colon U \to V$ between an open ball $U$ containing $x$ and an open ball $V\subset \matH^n$. 
We show that $D$ extends (uniquely) to an isometry $D\colon M\to\matH^n$. 

For every $y \in M$, choose an arc $\alpha\colon [0,1]\to M$ from $x$ to $y$. By compactness there is a partition $0=t_0 < t_1 < \ldots < t_k =1$ and for each $i=0,\ldots, k-1$ an isometry $D_i\colon U_i \to V_i$ from an open ball $U_i$ in $M$ containing $\alpha([t_i,t_{i+1}])$ to an open ball $V_i\subset \matH^n$. 

We may suppose that $U_0 \subset U$ and $D_0 = D|_{U_0}$. Inductively on $i$, we now modify $D_i$ so that $D_{i-1}$ and $D_i$ coincide on the component $C$ of $U_{i-1}\cap U_i$ containing $\alpha(t_i)$. To do so, note that
$$D_{i-1}\circ D_i^{-1}\colon D_i(C) \longrightarrow D_{i-1}(C)$$
is an isometry of open connected sets in $\matH^n$ and hence extends to an isometry of $\matH^n$. Then it makes sense to compose $D_i$ with $D_{i-1}\circ D_i^{-1}$, so that the new maps $D_{i-1}$ and $D_i$ coincide on $C$. Finally, we define $D(y) = D_{k-1}(y)$.

The proof that $D(y)$ is well-defined is a standard argument. First, it is easy to check that different partitions $0=t_0 <\ldots <t_k=1$ do not vary $D(y)$, just by considering a common refinement. Then we consider another path $\beta$ connecting $x$ to $y$. Since $M$ is simply-connected, there is a homotopy connecting $\alpha$ and $\beta$. The image of the homotopy is compact and is hence covered by finitely many open balls $U_i$ isometric to open balls $V_i\subset \matH^n$ via some maps $D_i$. By the Lebesgue number theorem, there is a $N>0$ such that in the grid in $[0,1]\times [0,1]$ of $\frac 1N \times \frac 1N$ squares, the image of every square is entirely contained in at least one $U_i$. We can now modify as above the isometries $D_i$ inductively on the grid, starting from the bottom-left square, so that they all glue up and show that $D(y)$ does not depend on $\alpha$ or $\beta$.

The resulting map $D\colon M \to \matH^n$ is a local isometry by construction. Since $M$ is complete, the map $D$ is a covering by Proposition \ref{local:isometry:prop}. Since $\matH^n$ is simply connected, the covering $D$ is a homeomorphism and $D$ is actually an isometry.
\end{proof}

The isometry $D\colon M \to \matH^n$ constructed in the proof is called a \emph{developing map}.
The same proof shows that every complete simply connected flat (or elliptic) $n$-manifold is isometric to $\matR^n$ (or $S^n$). \index{developing map}

\subsection{Complete hyperbolic manifolds}
We have determined the unique complete simply connected hyperbolic $n$-manifold, and we now look at complete hyperbolic manifolds with arbitrary fundamental group. We first note that if $\Gamma < \Iso(\matH^n)$ is a group of isometries that acts freely and properly discontinuously on $\matH^n$, the quotient manifold $\matH^n/_\Gamma$ has a natural Riemannian structure that promotes the covering
$$\pi\colon\matH^n \longrightarrow \matH^n/_\Gamma$$
to a local isometry, see Proposition \ref{quotient:prop}. The quotient $\matH^n/_\Gamma$ is a complete hyperbolic manifold. 
We now show that every complete hyperbolic manifold is realised in this way:

\begin{prop} \label{corto:prop}
Every complete hyperbolic $n$-manifold $M$ is isometric to $\matH^n/_\Gamma$ for some subgroup $\Gamma<\Iso(\matH^n)$ acting freely and properly discontinuously.
\end{prop}
\begin{proof}
The universal cover of $M$ inherits a Riemannian structure that is complete (by Proposition \ref{local:isometry:prop}), hyperbolic, and simply connected: hence it is isometric to $\matH^n$ by Theorem \ref{simply:teo}. The deck transformations $\Gamma$ of the covering $\matH^n\to M$ are necessarily locally isometries, hence isometries. We conclude that $M=\matH^n/_\Gamma$ and $\Gamma$ acts freely and properly discontinuously using Proposition \ref{equivalent:prop}.
\end{proof}

Note that $\Gamma$ is isomorphic to the fundamental group $\pi_1(M)$.

\begin{oss}
A group $\Gamma<\Iso(\matH^n)$ acts freely if and only if it does not contain elliptic isometries: that is, every non-trivial isometry in $\Gamma$ is either hyperbolic or parabolic.
\end{oss}
Note also that $\Gamma$ acts properly discontinuously if and only if it is discrete, see Proposition \ref{iff:discrete:prop}. 

\begin{oss}
The same proofs show that every complete flat or spherical $n$-manifold is isometric to $\matR^n/_\Gamma$ or $S^n/_\Gamma$ for some discrete group $\Gamma$ of isometries acting freely on $\matR^n$ or $S^n$. 
\end{oss}

\begin{cor} There is a natural 1-1 correspondence
$$\left\{\!
\begin{array}{c}{\rm complete\ hyperbolic} \\ {\rm manifolds\ } M {\rm\ up\ to\ isometry}\end{array}\!\right\} \longleftrightarrow
\left\{\!
\begin{array}{c}{\rm discrete\ subgroups\ } \Gamma < \Iso(\matH^n) \\ {\rm without\ elliptics} \\ {\rm up\ to\ conjugation}\end{array}\!\right\}$$
\end{cor}
\begin{proof}
When passing from the complete hyperbolic manifold $M$ to the group $\Gamma$, the only choice we made is an isometry between the universal cover of $M$ and $\matH^n$. Different choices produce conjugate groups $\Gamma$.
\end{proof}

\subsection{Discrete groups}
We investigate some basic properties of discrete groups $\Gamma$ of isometries of $\matH^n$.

\begin{ex} If $\Gamma<\Iso(\matH^n)$ is discrete then it is countable.
\end{ex}

Note that $\Gamma$ is not necessarily finitely generated. We denote by $\Gamma_p < \Gamma$ the stabiliser of a point $p\in \matH^n$.

\begin{prop} \label{discrete:prop}
Let $\Gamma <\Iso(\matH^n)$ be discrete and $p\in\matH^n$ be a point. The stabiliser $\Gamma_p$ is finite and the orbit $\Gamma(p) = \{g(p)\ |\ g\in \Gamma\}$ is discrete.
\end{prop}
\begin{proof}
Both are obvious consequence of the fact that $\Gamma$
acts properly discontinuously.
\end{proof}

Of course $\Gamma$ acts freely on $\matH^n$ if and only if $\Gamma_p$ is trivial for all $p\in\matH^n$. A set of subspaces in $\matH^n$ is \emph{locally finite} if every compact subset in $\matH^n$ intersects only finitely many of them.

\begin{prop} \label{dense:prop}
Let $\Gamma<\Iso(\matH^n)$ be discrete. The points $p\in\matH^n$ with trivial stabiliser $\Gamma_p$ form an open dense set in $\matH^n$.
\end{prop}
\begin{proof}
The fixed-points set $\Fix(g)$ of a non-trivial isometry $g$ is a proper subspace of $\matH^n$. The subspaces $\Fix(g)$ with $g\in\Gamma$ are locally finite: if infinitely many of them intersect a compact set they accumulate and $\Gamma$ does not act properly discontinuously. The complement of a locally finite set of proper subspaces is open and dense.
\end{proof}

Recall that a group \emph{has no torsion} if every non-trivial element has infinite order.

\begin{prop} \label{no:torsion:prop}
A discrete group $\Gamma < \Iso(\matH^n)$ acts freely on $\matH^n$ if and only if it has no torsion.
\end{prop}
\begin{proof}
By Proposition \ref{iterate:prop} parabolic and hyperbolic elements have infinite order. On the other hand, an elliptic element $g\in\Gamma$ has finite order since $\Gamma_p$ is finite for $p\in\Fix(g)$.
\end{proof}

\begin{cor}
The fundamental group of a complete hyperbolic manifold has no torsion.
\end{cor}

It is now time to exhibit some examples. Recall that $H^2$ is the half-plane model, see Section \ref{isometry:2:3:section}.

\begin{example}
The \emph{modular group} 
$$\Gamma = \PSLZ < \PSLR = \Iso^+(H^2)$$ 
consists of all matrices in $\PSLR$ having integer entries and is clearly a discrete subgroup. It does not act freely on $H^2$, however: the matrix $\matr 01{-1}0$ represents the elliptic transformation $z\mapsto -\frac 1z$ with fixed point $i$. \index{modular group}
\end{example} 

\subsection{Coverings}
We now make a simple but crucial observation: if $\Gamma < \Iso(\matH^n)$ acts freely and properly discontinuously, then also every subgroup $\Gamma' < \Gamma$ does; we get a manifolds covering
$$\matH^n/_{\Gamma'} \longrightarrow \matH^n/_\Gamma$$
whose degree $d$ is precisely the index of $\Gamma'$ in $\Gamma$.  Recall from Proposition \ref{volume:covering:prop} that we get
$$\Vol\big(\matH^n/_{\Gamma'}\big) = d \cdot \Vol\big(\matH^n/_{\Gamma}\big)$$
where some of the terms in the formula may be infinite. Moreover, every covering of a hyperbolic complete $M= \matH^n/_\Gamma$ is constructed in this way: there is a nice bijective correspondence
$$\big\{{\rm coverings\ of\ }M \big\} \longleftrightarrow
\big\{{\rm subgroups\ of\ }\Gamma \big\} .$$
This holds of course also for flat and spherical manifolds.

\subsection{Congruence subgroups}
We can now exhibit a family of two-dimensional hyperbolic manifolds. 

Pick an integer $m\geqslant 2$. Let $\SL_2(\matZ/_{m\matZ})$ be the group of $2\times 2$ matrices with coefficients in $\matZ/_{m\matZ}$ and determinant 1. We define the quotient 
$$\matP\SL_2(\matZ/_{m\matZ}) = \SL_2(\matZ/_{m\matZ})/_{\{ \pm I\}}.$$
The reduction modulo $m$ homomorphism $\matZ \to \matZ/_{m\matZ}$ induces the group homomorphisms $\SL_2(\matZ) \to \SL_2(\matZ/_{m\matZ})$ and 
$$\PSLZ \longrightarrow \matP\SL_2(\matZ/_{m\matZ}).$$
The kernel of this homomorphism is the \emph{principal congruence subgroup} $\Gamma(m)$ of $\PSLZ$. It is clearly discrete, since $\PSLZ$ is. It has finite index in $\PSLZ$ because $\matP\SL_2(\matZ/_{m\matZ})$ is finite. 
\index{principal congruence subgroup}
\begin{prop}
If $m\geqslant 4$ the group $\Gamma(m)$ acts freely on $H^2$.
\end{prop}
\begin{proof}
An element $A\in\Gamma(m)$ is a matrix $\matr abcd$ congruent to $\matr 1001$ modulo $m$. In particular $a+d$ is congruent to 2 modulo $m$, and hence is not $-1, 0, 1$. Therefore $A$ is never elliptic. 
\end{proof}

The quotient $\matH^2/_{\Gamma(m)}$ is a hyperbolic surface. We will construct many hyperbolic surfaces in Section \ref{surface:geometrisation:section} via some more geometric methods.

\subsection{Selberg's lemma}
The discrete group $\Gamma = \PSLZ$ does not act freely on $H^2$, but its finite-index normal subgroup $\Gamma(m)$ does as soon as $m\geqslant 4$. Is this the instance of a more general principle? Yes, it is. \index{Selberg lemma}

\begin{prop} \label{Selberg:prop}
Every finitely generated discrete group $\Gamma < \Iso(\matH^n)$ contains a finite-index normal subgroup $\Gamma'\triangleleft \Gamma$ that acts freely on $\matH^n$.
\end{prop}
\begin{proof}
The group $\Iso(\matH^n)$ is isomorphic to $O^+(n,1) < \GL(n+1,\matC)$, so Selberg's Lemma \ref{Selberg:lemma} applies to $\Gamma$, and it furnishes a finite-index torsion-free normal subgroup $\Gamma'\triangleleft \Gamma$. This subgroup acts freely by Proposition \ref{no:torsion:prop}.
\end{proof}

Every finitely generated discrete group $\Gamma < \Iso(\matH^n)$ contains at least one finite-index subgroup that acts freely. But how many such subgroups does it contain? Quite a lot, in fact.

\begin{prop} \label{FG:RF:prop}
Every finitely generated discrete group $\Gamma<\Iso(\matH^n)$ is residually finite.
\end{prop}
\begin{proof}
See Lemma \ref{RF:lemma}.
\end{proof}

The following corollary shows that there is an abundance of torsion-free subgroups. Although algebraic in nature, it has some remarkable geometric consequences, that will be revealed soon in Section \ref{finite:covers:subsection}.

\begin{cor}
Let $\Gamma < \Iso(\matH^n)$ be discrete and finitely generated. For every non-trivial $g \in \Gamma$ there is a finite-index normal subgroup $\Gamma'\triangleleft \Gamma$ that acts freely on $\matH^n$ and does not contain $g$.
\end{cor}

\section{Polyhedra} 
A polyhedron in $\matH^n$ is a natural geometric object, that may be used to visualise discrete groups in $\Iso(\matH^n)$ and hence hyperbolic manifolds. 
Polyhedra may sometimes be combined to form some tessellations of the space. 

\subsection{Polyhedra} \label{polyhedra:subsection}
A \emph{half-space} in $\matH^n$ is the closure of one of the two portions of space delimited by a hyperplane. We say that a set of half-spaces is locally finite if their boundary hyperplanes are. \index{polyhedron}

\begin{defn}
A $n$-dimensional \emph{polyhedron} $P$ in $\matH^n$ is the intersection of a locally finite set of half-spaces. We also assume that $P$ has non-empty interior.
\end{defn}

A subset $S\subset \overline{\matH^n}$ is \emph{convex} if $x,y\in S$ implies that the segment connecting $x,y$ is also contained in $S$ (such a segment is a half-line or a line if one or both points lie in $\partial \matH^n$). Every polyhedron $P$ is clearly convex because it is the intersection of convex sets. Its closure $\overline P$ in $\overline{\matH^n}$ is also convex.
\index{convex subset}

Let $H\subset \matH^n$ be a half-space containing the polyhedron $P$. If non-empty, the intersection $F=\partial H\cap P$ is called a \emph{face} of $P$. The \emph{supporting subspace} of $F$ is the smallest subspace of $\matH^n$ containing $F$; the \emph{dimension} of a face is the dimension of its supporting subspace. A face of dimension $0$, $1$, and $n-1$ is called a \emph{vertex}, an \emph{edge}, and a \emph{facet}.

\begin{ex} \label{intersection:faces:ex}
If non-empty, the intersection of faces of $P$ is a face.
\end{ex}

\begin{ex} Every $k$-dimensional face is a polyhedron in its supporting $k$-dimensional space.
\end{ex}

The \emph{convex hull} of a set $S\subset \overline{\matH^n}$ is the intersection of all the convex sets containing $S$. \index{convex hull}

\begin{ex} \label{hull:ex}
The convex hull of finitely many points in $\matH^n$ that are not contained in a hyperplane is a compact polyhedron. Conversely, every compact polyhedron has finitely many vertices and is the convex hull of them.
\end{ex}

Everything we said holds with no modifications for $\matR^n$. On $S^n$ some care should be taken: some definitions need to be modified slightly to take into account the annoying presence of antipodal points. We gloss over this technical point.

\subsection{Finite polyhedra}
We now enlarge slightly the class of compact polyhedra by admitting finitely many vertices at infinity. \index{polyhedron!finite polyhedron}

\begin{defn}
A \emph{finite polyhedron} is the convex hull of finitely many points $x_1,\ldots, x_k \in \overline{\matH^n}$ that are not contained in the closure of a hyperplane.
\end{defn}

The $x_i$'s that lie in $\partial \matH^n$ are called \emph{ideal vertices}, while the usual vertices of $P$ are the \emph{finite} or \emph{actual vertices}. The ideal vertices form the set $\overline P \setminus P$.

\begin{ex}
A finite polyhedron $P$ has finitely many faces and is the convex hull of its ideal and finite vertices.
\end{ex}

We want to estimate the volume of finite polyhedra, and to do this we need a lemma. Given a horosphere $O$ centred at $p\in\partial \matH^n$ and a domain $D\subset O$, the \emph{cone} $C$ of $D$ over $p$ is the union of all the half-lines exiting from $D$ and pointing towards $p$, see Figure \ref{volume:fig}.

\begin{lemma} \label{O:lemma}
Let $O$ be a horosphere centred at $p\in \partial \matH^n$, $D\subset O$ any domain and $C$ the cone over $D$. The following equality holds:
$$\Vol(C) = \frac {\Vol_O (D)}{n-1}$$
where $\Vol_O$ is the $(n-1)$-dimensional volume in the flat $(n-1)$-manifold $O$.
\end{lemma}
\begin{proof}
Let $O$ have some height $x_n=h$ as in Figure \ref{volume:fig}-(left). We obtain
$$\Vol(C) = \int_D dx \int_{h}^\infty \frac 1{t^{n}}dt = \frac 1{n-1}\int_D \frac{dx}{h^{n-1}} = \frac 1{n-1} \cdot \Vol_O(D).$$
The proof is complete.
\end{proof}

\begin{figure}
\begin{center}
\includegraphics[width = 12 cm] {\iftoggle{BW}{volume-BW}{volume}} 
\nota{The cone $C$ over a domain $D\subset O$ has volume proportional to the area of $D$ (left). If the domain is compact, the cone has finite volume: therefore a finite polyhedron has finite volume (right).}
\label{volume:fig}
\end{center}
\end{figure}

We now turn to finite polyhedra.

\begin{prop}
Every finite polyhedron has finite volume.
\end{prop}
\begin{proof}
For every ideal vertex of $P$, a small horoball centred at $p$ intersects $P$ into a cone that has finite volume. The polyhedron $P$ decomposes into finitely many cones and a bounded region, see Figure \ref{volume:fig}-(right).
\end{proof}

A finite polyhedron without finite vertices is called an \emph{ideal polyhedron}.
\index{polyhedron!ideal polyhedron}

\subsection{Polygons}
A polygon is just a polyhedron of dimension two. In contrast with Euclidean geometry, a strikingly simple formula relates the area of a finite polygon with its inner angles. We define the inner angle of an ideal vertex to be zero. 

\begin{figure}
\begin{center}
\includegraphics[width = 10 cm] {\iftoggle{BW}{triangle_area-BW}{triangle_area}}
\nota{A triangle with at least an ideal vertex (left). The area of a triangle with finite vertices can be derived as the area difference of triangles with one ideal vertex (right).}
\label{triangle_area:fig}
\end{center}
\end{figure}

\begin{prop} \label{polygon:area:prop}
A polygon $P$ with inner angles $\alpha_1,\ldots, \alpha_n$ has area
$$\Area (P) = (n-2)\pi - \sum_{i=1}^n \alpha_i.$$
\end{prop}
\begin{proof}
Every polygon decomposes into triangles, and it suffices to prove the formula on these. Consider first a triangle $T$ with at least one vertex at infinity. We use the half-plane model and send this vertex to $\infty$ as in Figure \ref{triangle_area:fig}-(left). We suppose that the \iftoggle{BW}{grey}{red} dot is the origin of $\matR^2$, hence 
$$T = \big\{(r\cos\theta, y)\ | \ \beta\leqslant \theta \leqslant \pi-\alpha, \ y\geqslant r\sin\theta\big\}$$
and we get
\begin{align*}
\Area(T) & = \int_T \frac 1{y^2}dxdy = \int_{\pi-\alpha}^{\beta} \int_{r\sin\theta}^\infty \frac{-r\sin \theta}{y^2}dyd\theta \\
 & = \int_{\pi-\alpha}^\beta -r\sin\theta \left[-\frac 1y\right]_{r\sin\theta}^\infty d\theta = 
\int_{\beta}^{\pi-\alpha} \frac{r\sin\theta}{r\sin\theta}d\theta \\
 & = \int_{\beta}^{\pi-\alpha} 1 = \pi-\alpha-\beta.
\end{align*} 
The area of a triangle with finite vertices $ABC$ is deduced as in Figure \ref{triangle_area:fig}-(right) using the formula
$$\Area (ABC) = \Area(AB\infty) + \Area(BC\infty) - \Area (AC\infty).$$
The proof is complete.
\end{proof}

The sum of the inner angles of a hyperbolic polygon is strictly smaller than that of a Euclidean polygon with the same number of sides, and the difference between these two numbers is precisely its area. 

\begin{cor} Every ideal triangle has area $\pi$.
\end{cor}

\subsection{Platonic solids} \label{platonici:subsection}
The theory of three-dimensional polyhedra in $\matH^3$ is very rich: we limit ourselves to the study of the platonic solids.

The five Euclidean platonic solids are the regular tetrahedron, the cube, the regular octahedron, icosahedron, and dodecahedron. We now see that each platonic solid $P\subset \matR^3$ generates a nice continuous family of solids in the three geometries $\matH^3, \matR^3$, and $S^3$.

To construct this family we fix any point $x$ in $\matH^3$ and represent $P$ centred at $x$ with varying size. To do this, we put $P$ inside the Euclidean tangent space $T_x\matH^3$ centred at the origin and with some radius $t>0$. Consider the image of its vertices by the exponential map and take their convex hull. We indicate by $P(-t)$ the resulting polyhedron in $\matH^3$. The polyhedron $P(-t)$ is combinatorially equivalent to $P$ and has the same symmetries of $P$.

We extend this family to the other geometries as follows. The polyhedron $P(0)$ is the Euclidean $P$ (unique up to dilations), and $P(t)$ with $t>0$ is the spherical $P$, constructed as above with $S^3$ instead of $\matH^3$: we fix $x\in S^3$, take a copy of $P$ inside $T_x S^3$ with radius $t$, project its vertices, and take the convex hull. We define the spherical $P(t)$ only for $t \in (0, \frac \pi 2]$: when $t= \frac\pi 2$ it degenerates to a hemisphere. We also define $P(-\infty)$ as the ideal hyperbolic platonic solid obtained by sending all the vertices at infinity.

\begin{figure}
\begin{center}
\includegraphics[width = 3.5 cm] {\iftoggle{BW}{Icosahedron-BW}{256px-Icosahedron}} \qquad
\includegraphics[width = 3.5 cm] {\iftoggle{BW}{Dodecahedron-BW}{245px-Dodecahedron}}
\nota{The regular icosahedron and dodecahedron.}
\label{icosaedro:fig}
\end{center}
\end{figure}

We have defined a polyhedron $P(t)$ for all $t\in [-\infty, \frac \pi 2]$, that lies in $\matH^3, \matR^3, S^3$ depending on whether $t$ is negative, null, or positive. In some sense the polyhedron $P(t)$ depends continuously on $t$ also when it crosses the value $t=0$: when $t\to 0$ the polyhedron in $\matH^3$ or $S^3$ shrinks, and every polyhedron tends to a Euclidean one when shrunk. 

In particular the dihedral angle $\theta(t)$ of $P(t)$ varies continuously with $t\in [-\infty, \frac \pi2]$. The function $\theta(t)$ is strictly monotone and we now determine its image. The \emph{vertex valence} of $P$ is the number of edges incident to each vertex.

\begin{prop} \label{dihedral:platonic:prop}
Let $n\in\{3,4,5\}$ be the vertex valence of $P$. It holds 
$$\theta \left(\left[-\infty, \frac \pi 2\right]\right) = \left[\frac{n-2}n \pi, \pi\right].$$
\end{prop}
\begin{proof}
Since $\theta$ is continuous and monotone increasing, it suffices to show that $\theta(-\infty) = \frac{n-2}n\pi$ and $\theta(\frac \pi 2) = \pi$.

By intersecting the ideal polyhedron $P(-\infty)$ with a small horosphere $O$ centred at an ideal vertex $v$ we get a regular $n$-gon in the Euclidean plane $O$, with interior angles $\frac{n-2}n\pi$. The dihedral angle at an edge $e$ incident to $v$ is measured by intersecting $P(-\infty)$ with any hypersurface orthogonal to $e$: since $O$ is orthogonal to $e$ we get $\theta(-\infty) =  \frac{n-2}n\pi$. 

The polyhedron $P(\frac \pi 2)$ is a hemisphere and hence $\theta(\frac \pi 2) = \pi$.
\end{proof}

For some values of $t$, the platonic solid $P(t)$ may have nice dihedral angles $\theta(t)$, for instance angles that divide $2\pi$. In the Euclidean world, the only platonic solid with such nice dihedral angles is the cube. In the hyperbolic and spherical world we find more.

The platonic solids with dihedral angles that divide $2\pi$ are listed in Table \ref{platonic:table}. The table is just a consequence of Proposition \ref{dihedral:platonic:prop}: it suffices to know the Euclidean dihedral angles for $P$, and all the angles bigger (smaller) than this value are spherical (hyperbolic).

\begin{table} 
\begin{center}
\begin{tabular}{c||cccc}
\phantom{\Big|} polyhedron & $\theta = \frac \pi 3$ & $\theta = \frac{2\pi}5$ & $\theta = \frac \pi 2$ &
$\theta = \frac{2\pi}3$ \\
\hline\hline
\rule{0pt}{3ex}
\phantom{\big|}  tetrahedron & ideal $\matH^3$ & $S^3$ & $S^3$ & $S^3$ \\
\phantom{\big|}  cube & ideal $\matH^3$ & $\matH^3$ & $\matR^3$ & $S^3$ \\
\phantom{\big|}  octahedron & & & ideal $\matH^3$ & $S^3$ \\
\phantom{\big|} icosahedron & & & & $\matH^3$ \\
\phantom{\big|} dodecahedron & ideal $\matH^3$ & $\matH^3$ & $\matH^3$ & $S^3$ \\
\end{tabular}
\end{center}
\nota{The platonic solids with dihedral angle $\theta$ that divide $2\pi$.}
\label{platonic:table}
\end{table}

In particular, there are four right-angled platonic solids: the spherical tetrahedron, the  Euclidean cube, the hyperbolic dodecahedron, and the ideal hyperbolic octahedron.

\section{Tessellations}
A tessellation is a nice paving of $\matH^n$ made of polyhedra. Not only tessellations are beautiful objects, but they are also tightly connected with discrete groups of $\Iso(\matH^n)$ and hence with hyperbolic manifolds.

%This is their definition. 
\index{tessellation}

\begin{defn}
A \emph{tessellation} of $\matH^n$ (or $\matR^n$, $S^n$) is a locally finite set of polyhedra that cover the space and may intersect only in common faces.
\end{defn}

\begin{figure}
\begin{center}
\includegraphics[width = 5 cm] {\iftoggle{BW}{Uniform_tiling_532-t012-BW}{Uniform_tiling_532-t012}} \quad
\includegraphics[width = 5 cm] {\iftoggle{BW}{Truncated_octahedra-BW}{Truncated_octahedra}}
\nota{A tessellation of $S^2$ into squares, hexagons, and decagons, and a tessellation of $\matR^3$ into truncated octahedra.}
\label{tiling:fig}
\end{center}
\end{figure}

Some examples in spherical and Euclidean space are shown in Figure \ref{tiling:fig}. We now construct some families explicitly.

\subsection{Triangles} \label{triangular:subsection}
We want to define some nice tessellations of $\matH^2, \matR^2$, and $S^2$ into triangles. The following exercise says that every triple of acute angles is realised by some triangle in the appropriate geometry.

\begin{ex}
Given three real numbers $0<\alpha, \beta, \gamma \leqslant \frac \pi 2$ there is a triangle $\Delta$ with inner angles $\alpha, \beta, \gamma$ inside $\matH^2, \matR^2$, or $S^2$ depending on whether the sum $\alpha + \beta + \gamma$ is smaller, equal, or bigger than $\pi$. 
\end{ex}

Let $a,b,c \geqslant 2$ be three natural numbers and $\Delta$ be a triangle with inner angles $\frac \pi a, \frac \pi b, \frac \pi c$. The triangle $\Delta$ lies in $\matH^2$, $\matR^2$, or $S^2$ depending on whether $\frac 1a+\frac 1b + \frac 1c$ is smaller, equal, or bigger than 1. In all cases, by mirroring iteratively $\Delta$ along its edges we construct a tessellation of the space. 


\begin{figure}
\begin{center}
\includegraphics[width = 3.5 cm] {\iftoggle{BW}{233_BW}{233}} \quad
\includegraphics[width = 3.5 cm] {\iftoggle{BW}{234_BW}{234}} \quad 
\includegraphics[width = 3.5 cm] {\iftoggle{BW}{235_BW}{235}} 
\nota{The tessellations $(2,3,3)$, $(2,3,4)$ and $(2,3,5)$ of the sphere.}
\label{triangular:fig}
\end{center}
\end{figure}

The triples realisable in $S^2$ are $(2,2,c)$, $(2,3,3)$, $(2,3,4)$, and $(2,3,5)$: the last three tessellations are shown in Figure \ref{triangular:fig} and are connected to the platonic solids. They consist of $24$, $48$, and $120$ triangles. 

\begin{figure}
\begin{center}
\includegraphics[width = 3.5 cm] {\iftoggle{BW}{Tile_V46b_BW}{200px-Tile_V46b}} \quad
\includegraphics[width = 3.5 cm] {\iftoggle{BW}{Tile_V488_bicolor_BW}{200px-Tile_V488_bicolor}} \quad 
\includegraphics[width = 3.5 cm] {\iftoggle{BW}{1024px-Tile_36_BW}{200px-Tile_36}} 
\nota{The tessellations $(2,3,6)$, $(2,4,4)$, and $(3,3,3)$ of the Euclidean plane.}
\label{triangular2:fig}
\end{center}
\end{figure}

The triples realisable in $\matR^2$ are $(2,3,6)$, $(2,4,4)$, and $(3,3,3)$: the tessellations are shown in Figure \ref{triangular2:fig}. There are infinitely many triples realisable in $\matH^2$, and some are shown in Figure \ref{triangular3:fig}. 

\begin{figure}
\begin{center}
\includegraphics[width = 3.7 cm] {\iftoggle{BW}{Order-3_heptakis_heptagonal_tiling-BW}{240px-Order-3_heptakis_heptagonal_tiling}} \quad
\includegraphics[width = 3.7 cm] {\iftoggle{BW}{Order-4_bisected_pentagonal_tiling-BW}{240px-Order-4_bisected_pentagonal_tiling}} \quad 
\includegraphics[width = 3.7 cm] {\iftoggle{BW}{Uniform_dual_tiling_433-t012-BW}{240px-Uniform_dual_tiling_433-t012}} 
\nota{The tessellations $(2,3,7)$, $(2,4,5)$, and $(3,3,4)$ of the hyperbolic plane.}
\label{triangular3:fig}
\end{center}
\end{figure}

In the hyperbolic plane we can also use triangles with vertices at infinity, which have inner angle zero by definition.
\begin{ex}
Given three real numbers $0\leqslant \alpha, \beta, \gamma \leqslant \frac \pi 2$ with sum smaller than $\pi$ there is a triangle $\Delta\subset \matH^2$ with inner angles $\alpha, \beta, \gamma$.
\end{ex}

For any triple $(a,b,c)$ of numbers in $\matN \cup \{\infty \}$ with $\frac 1a + \frac 1b + \frac 1c < 1$ we may take the triangle $\Delta \subset \matH^2$ with inner angles $\frac \pi a$, $\frac \pi b$, $\frac \pi c$ and reflect it iteratively to get a tessellation of $\matH^2$. 

The triple $(\infty, \infty, \infty)$ gives a nice tessellation into ideal triangles called the \emph{Farey tessellation} and shown in Figure \ref{tiling_Farey:fig}. 
\index{tessellation!Farey tessellation}
 
\begin{figure}
\begin{center}
\includegraphics[width = 5 cm] {H2checkers_iii.png}
\nota{The Farey tessellation $(\infty, \infty, \infty)$ of the hyperbolic plane.}
\label{tiling_Farey:fig}
\end{center}
\end{figure}

\subsection{Platonic solids}
We now turn to tessellations of 3-dimensional spaces. Table \ref{platonic:table} displays a finite list of platonic solids $P$ with dihedral angles $\frac{2\pi}k$ contained in $\matH^3$, $\matR^3$, or $S^3$. For each solid in the list, by reflecting iteratively $P$ along its faces we get a tessellation of the space. 

The Figures \ref{honeycomb_dodecahedron_534:fig}, \ref{honeycomb_dodecahedron_536:fig}, and \ref{honeycomb_dodecahedron_533:fig} show the tessellations of $\matH^3, \matH^3,$ and $S^3$ by regular dodecahedra with dihedral angles $\frac \pi 2$, $\frac \pi 3$, and $\frac{2\pi}3$. The first two tessellations contain infinitely many polyhedra and appear in the figures as seen by an observer floating inside $\matH^3$. The reader is invited to observe the differences between the first two tessellations of the hyperbolic space. The third tessellation contains finitely many polyhedra and the figure shows its stereographic projection in $\matR^3$.

\begin{figure}
\begin{center}
\includegraphics[width = 9 cm] {\iftoggle{BW}{H3_534_CC_center-BW}{H3_534_CC_center}}
\nota{The tessellation of $\matH^3$ into right-angled regular dodecahedra, seen from inside $\matH^3$.}
\label{honeycomb_dodecahedron_534:fig}
\end{center}
\end{figure}

\begin{figure}
\begin{center}
\includegraphics[width = 9 cm] {\iftoggle{BW}{H3_536_CC_center-BW}{H3_536_CC_center}}
\nota{The tessellation of $\matH^3$ into ideal regular dodecahedra with dihedral angle $\frac \pi 3$, seen from inside $\matH^3$. Note that all vertices lie at infinity.}
\label{honeycomb_dodecahedron_536:fig}
\end{center}
\end{figure}

\begin{figure}
\begin{center}
\includegraphics[width = 9 cm] {\iftoggle{BW}{Stereographic_polytope_120cell_faces-BW}{600px-Stereographic_polytope_120cell_faces}}
\nota{The tessellation of $S^3$ into 120 regular dodecahedra with dihedral angle $\frac {2\pi} 3$, transposed into $\matR^3$ via the stereographic projection, that transforms straight faces into round ones but preserves the angles.}
\label{honeycomb_dodecahedron_533:fig}
\end{center}
\end{figure}

\subsection{Regular tessellations} \label{regular:subsection}
Platonic solids are \emph{regular}, that is they have many symmetries. A notion of regularity may be defined in all dimensions for polyhedra and tessellations as follows.
\index{tessellation!regular tessellation} \index{polyhedron!regular polyhedron}

Every tessellation in $\matH^n$, $\matR^n$, or $S^n$ has a symmetry group, consisting of all the isometries of the ambient space that preserve it. A \emph{flag} is a sequence of faces $f_0 \subset \ldots \subset f_n$ of the tessellation with $\dim f_i = i$. A tessellation is \emph{regular} if the symmetry group acts transitively on flags. 

These definitions apply also to polyhedra. We note that a regular polyhedron in $\matR^n$, considered up to similarities, may be interpreted as a regular tessellation in $S^{n-1}$, and vice versa.\index{Schl\"afli notation}

The \emph{Schl\"afli notation} beautifully encodes various regular tessellations, in all dimensions and in all geometries. The symbol $\{n\}$ with $n\geqslant 3$ denotes the regular $n$-gon, that is a regular tessellation of $S^1$ by $n$ congruent arcs, and by extension $\{\infty \}$ denotes the regular tessellation of $\matR$ by infinitely many congruent segments.

The symbol $\{p,q\}$ denotes a regular tessellation in $\matH^2, \matR^2,$ or $S^2$ by $p$-gons where $q$ of them meet at every vertex. The five platonic solids are
$$\{3,3\}, \{3,4\}, \{3,5\}, \{4,3\}, \{5,3\},$$
the symbols $\{3,6\}$, $\{4,4\}$, and $\{6,3\}$ denote the two regular tessellations of $\matR^2$ into equilateral triangles, squares, and regular hexagons, and every other pair $\{p,q\}$ denotes a tessellation of regular $p$-gons in $\matH^2$ with angles $\frac {2\pi}q$, that meet at $q$ at every vertex. We can also interpret $\{p, \infty\}$ as a regular tessellation of $\matH^2$ into ideal regular $p$-gons (the Farey tessellation is $\{3,\infty\}$) and $\{\infty, q\}$ as a regular tessellation of $\infty$-gons that meet at $q$ at every vertex, whose edges form a $q$-regular tree (a tree where $q$ edges meet at every vertex). Finally $\{\infty, \infty\}$ is a regular tessellation of ideal $\infty$-gons. See Figure \ref{infty:tessellations:fig}

\begin{figure}
\begin{center}
\includegraphics[width = 3.8 cm] {\iftoggle{BW}{600px-H2chess_iiif}{600px-H2chess_iiif}} \quad
\includegraphics[width = 3.8 cm] {\iftoggle{BW}{600px-H2chess_24ic}{600px-H2chess_24ic}} \quad 
\includegraphics[width = 3.8 cm] {\iftoggle{BW}{600px-H2chess_2iib}{600px-H2chess_2iib}} 
\nota{The tessellations of $\matH^2$ with Schl\"afli symbols $\{4,\infty\}$, $\{\infty, 4\}$, and $\{\infty, \infty\}$.}
\label{infty:tessellations:fig}
\end{center}
\end{figure}

The symbol $\{p,q,r\}$ denotes a regular tessellation in $\matH^3, \matR^3,$ or $S^3$ by polyhedra $\{p,q\}$ where $r$ of them meet at every edge. We deduce from Table \ref{platonic:table} that the regular tessellations of $S^3$ are
$$\{3,3,3\}, \{3,3,4\}, \{3,3,5\}, \{3,4,3\}, \{4,3,3\}, \{5,3,3\}.$$
These are the six regular polyhedra (usually called \emph{polytopes}) in dimension four: we will encounter them again in Chapter \ref{regular:polytopes:subsection}. The symbol $\{4,3,4\}$ indicates the tessellation of $\matR^3$ into cubes, and
$$\{3,3,6\}, \{3,4,4\}, \{3,5,3\}, \{4,3,5\}, \{4,3,6\}, \{5,3,4\}, \{5,3,5\}, \{5,3,6\}$$
denote the tessellations of $\matH^3$ into platonic solids. It is also possible to interpret more triples $\{p,q,r\}$ in an appropriate way, sometimes by representing the vertices of the polyhedra as space-like vectors in the hyperboloid model.

Every regular tessellation has a \emph{dual} regular tessellation obtained by taking the barycenters of all the cells involved. The dual of $\{a,b, \ldots, z\}$ is $\{z, \ldots, b,a\}$. 
Duals of ideal tessellations involve infinite polyhedra: for instance $\{6,3,3\}$ is shown in Figure \ref{honeycomb_dodecahedron_633:fig}.

\begin{figure}
\begin{center}
\includegraphics[width = 9 cm] {\iftoggle{BW}{H3_633_FC_boundary-BW}{H3_633_FC_boundary}}
\nota{In the tessellation $\{6,3,3\}$ of $\matH^3$ every polyhedron has infinitely many hexagonal faces, whose vertices all lie in a single horosphere.}
\label{honeycomb_dodecahedron_633:fig}
\end{center}
\end{figure}

The story continues in four dimensions: we denote by $\{p,q,r,s\}$ a regular tessellation made of polytopes of type $\{p,q,r\}$ that meet in $s$ at every codimension-two face. With the same techniques of Section \ref{platonici:subsection} we can identify the regular polytopes in $\matH^4, \matR^4$ and $S^4$ with dihedral angles that divide $2\pi$, and the corresponding tessellations of the ambient space -- that we are unfortunately unable to see.

\subsection{Voronoi tessellations} \label{Voronoi:subsection}
In the previous section we have constructed some tessellations by exploiting the symmetries of regular polyhedra: we can probably extend these methods to less symmetric polyhedra, but how far can we go? How can we construct highly non-regular tessellations?

There is a strikingly simple procedure that transforms every discrete set $S$ of points in $\matH^n$ into a tessellation, called the \emph{Voronoi tessellation} of $S$. The construction goes as follows. \index{tessellation!Voronoi tessellation}

For every point $p \in S$ we define
$$D(p) = \big\{ q\in \matH^n\ \big|\ d(q,p)\leqslant d(q,p') \quad \forall p'\in S \big\}.$$

\begin{prop}
The set $D(p)$ is a polyhedron and the polyhedra $D(p)$ as $p\in S$ varies form a tessellation of $\matH^n$.
\end{prop}
\begin{proof}
It is an easy exercise to show that the points in $\matH^n$ having the same distance from two distinct fixed points form a hyperplane. For every $p'\in S$ different from $p$ we define the half-space
$$H_{p'} = \big\{q \in \matH^n\ \big| \ d(q,p) \leqslant d(q,p') \big\}.$$
The set $D(p)$ is the intersection of the half-paces $H_{p'}$ as $p'$ varies in $S\setminus \{p\}$.

Since $S$ is discrete, these half-spaces are locally finite (there are finitely many points in $S$ at bounded distance from $p$, hence finitely many hyperplanes). Therefore $D(p)$ is a polyhedron. Since $S$ is discrete, every point $q\in \matH^n$ has at least one nearest point $p\in S$: therefore the polyhedra $D(p)$ cover $\matH^n$ as $p\in S$ varies. 

It remains to prove that the polyhedra $D(p)$ intersect along common faces. We have $D(p) \cap D(p') = D(p) \cap \partial H_{p'}$ and hence $D(p) \cap D(p')$ is either empty or a face of $D(p)$. The case of multiple intersections follows from Exercise \ref{intersection:faces:ex}.
\end{proof}

\begin{figure}
\begin{center}
\includegraphics[width = 5 cm] {\iftoggle{BW}{Coloured_Voronoi_2D-BW}{Coloured_Voronoi_2D}} 
\nota{A Voronoi tessellation of the Euclidean plane.}
\label{Voronoi:fig}
\end{center}
\end{figure}

Voronoi tessellations of course make sense also in $\matR^n$ and $S^n$, see Figure \ref{Voronoi:fig}. We have just proved that tessellations are not exoteric, but quite ordinary objects, and we now show that they are useful to study hyperbolic manifolds.

\section{Fundamental domains}
We turn back to our discrete subgroups $\Gamma < \Iso(\matH^n)$. The geometry of a discrete $\Gamma$ is nicely controlled by some polyhedra, called \emph{fundamental domains}.
We introduce these objects and make some important examples.

\subsection{Fundamental and Dirichlet domains}
Let $\Gamma$ be a discrete group of isometries of $\matH^n$. The group $\Gamma$ may or may not act freely on $\matH^n$.  \index{fundamental domain}
\begin{defn}
A \emph{fundamental domain} for $\Gamma$ is a polyhedron $D \subset\matH^n$ whose translates $g(D)$ as $g\in \Gamma$ varies are distinct and form a tessellation of the space $\matH^n$. 
\end{defn}
If $D$ is a fundamental domain, the group $\Gamma$ acts freely and transitively on the tessellation $\big\{g(D)\big\}_{g \in \Gamma}$, so in particular the polyhedra $g(D)$ are all isometric. 

We describe a procedure that builds a fundamental domain for any discrete subgroup $\Gamma < \Iso(\matH^n)$. Pick  a point $p\in \matH^n$ 
with trivial stabiliser $\Gamma_p$, which exists by Proposition \ref{dense:prop}. The group $\Gamma$ acts freely and transitively on the orbit $\Gamma(p)$, which is discrete by Proposition \ref{discrete:prop}.

The discrete orbit $\Gamma(p)$ defines a $\Gamma$-invariant Voronoi tessellation of $\matH^n$, and every polyhedron of the tessellation is a fundamental domain. The polyhedron $D(p)$ of the tessellation containing $p$ is called the \emph{Dirichlet domain} for $\Gamma$ centred at $p$. By construction we have \index{Dirichlet domain}
$$D(g(p)) = g(D(p))$$
for all $g\in \Gamma$. We have proved in particular the following.

\begin{prop}
Every discrete group $\Gamma < \Iso(\matH^n)$ has a fundamental domain.
\end{prop}

Fundamental domains are far from being unique. The Dirichlet domain $D(p)$ depends on $p$ in a continuous fashion, and many fundamental domains are not Dirichlet domains.

\begin{figure}
\begin{center}
\includegraphics[width = 11 cm] {\iftoggle{BW}{ModularGroup-FundamentalDomain-01-BW}{ModularGroup-FundamentalDomain-01}}
\nota{The shadowed triangle (with one ideal vertex) is a fundamental domain for $\Gamma = \PSLZ$ acting on the half-space $H^2$. The translates of the fundamental domain form the tessellation shown.}
\label{modular_group:fig}
\end{center}
\end{figure}

\begin{ex}
Prove that the shadowed triangle in Figure \ref{modular_group:fig} is a fundamental domain for the modular group $\PSLZ$. 
\end{ex}

Everything we say also holds for $\matR^n$ and $S^n$.

\subsection{Fundamental domain of manifolds}
Let $\Gamma< \Iso(\matH^n)$ be a discrete subgroup that acts freely on $\matH^n$ and $D$ be a fundamental domain for $\Gamma$. We can get some information on the hyperbolic manifold $M=\matH^n/_\Gamma$ by looking at $D$. 

\begin{prop}
Let $M=\matH^n/_\Gamma$ be a hyperbolic manifold and $D$ a fundamental domain for $\Gamma$. The projection $\pi\colon \matH^n \to M$ restricts to a surjective map $D \to M$ that sends $\interior D$ isometrically onto an open dense subset of $M$. In particular we have
$$\Vol(D) = \Vol (M).$$
If $D$ is a Dirichlet domain, it is compact if and only if $M$ is.
\end{prop}
\begin{proof}
The translates $g(D)$ cover $\matH^n$, hence the projection $D \to M$ is surjective.
The translates $g(\interior D)$ are disjoint, hence $\interior D$ is sent isometrically inside $M$. The boundary $\partial D$ has measure zero, hence $\Vol(D) = \Vol(\interior D) = \Vol(M)$. 

If $D$ is compact then $M$ clearly is. If $D = D(p)$ is a Dirichlet domain and $M$ is compact, $M$ has finite diameter $\delta$ and hence every point in $\matH^n$ is at distance $\leqslant \delta$ of some point in the orbit of $p$. Hence $D(p)$ is contained in the closure of $B(p,\delta)$.
\end{proof}

A fundamental domain $D$ alone however does not determine $M$: we will soon see in Section \ref{piatte:subsection} that a square in $\matR^2$ is the fundamental domain of two non homeomorphic compact flat surfaces.

\subsection{Asteroids}
The facets of a fundamental domain $D$ are naturally partitioned into isometric pairs, as follows. Ever facet $f$ of $D$ is incident to $D$ and to another fundamental domain $g(D)$ of the tessellation. The isometry $g^{-1}$ sends $g(D)$ to $D$ and hence sends $f$ to another facet $f'$ of $D$. One checks immediately that $(f')' = f$, so $f$ and $f'$ are paired isometrically. 

One should think of $M$ as obtained from $D$ by identifying these facets in pairs: we can picture an observer -- say, a spaceship -- floating and traveling inside $M$ by visualising it in $D$, and jumping from $f$ to $f'$ every time it crosses the interior of a facet $f$, like in the 1979 video game \emph{Asteroids}.

\subsection{Spine and cut locus}
Another picture that may help to understand $M$ geometrically is the \emph{spine}\index{spine} $S \subset M$ defined as $S=\pi(\partial D)$. This is a $(n-1)$-dimensional object in $M$ whose complement is an open ball. When $D = D(p)$ is the Dirichlet domain of a point $p$, the spine $S$ is called the \emph{cut locus} of the point $q=\pi(p) \in M$. \index{cut locus}

\begin{ex}
The cut locus of $q \in M$ consists of all points $q'$ such that there are at least two geodesics of minimal length connecting $q$ to $q'$.
\end{ex}

We now exhibit some important examples of discrete groups of isometries in the three geometries $\matH^n$, $\matR^n$, and $S^n$, and study their fundamental domains. 

\subsection{Triangle groups} \label{triangle:groups:subsection}
Let $a,b,c \geqslant 2$ be natural numbers and $\Delta$ be the triangle in $\matH^2, S^2$, or $\matR^2$ with inner angles $\frac \pi a$, $\frac \pi b$, $\frac \pi c$. By reflecting iteratively $\Delta$ along its sides we get a tessellation $T$, see Section \ref{triangular:subsection}. \index{group!triangle group}

The \emph{triangle group} $\Gamma = \Gamma (a,b,c)$ is the group of isometries of $S^2$, $\matR^2$, or $\matH^2$ generated by the reflections $x,y,z$ along the three sides of $\Delta$ opposite to the vertices with inner angles $a,b,c$, respectively.

\begin{prop} \label{triangle:gamma:prop}
The triangle group $\Gamma (a,b,c)$ acts freely and transitively on the triangles of the tessellation $T$. Hence it is discrete and $\Delta$ is a fundamental domain for $\Gamma$. A presentation for the group is
$$\langle x, y,z\ |\ x^2, y^2, z^2, (xy)^c, (yz)^a, (zx)^b \rangle.$$
\end{prop}
\begin{proof}
We restrict for simplicity to the hyperbolic case, the others being analogous. It is convenient to construct the tessellation $T$ abstractly (this also furnishes a rigorous proof that by mirroring $\Delta$ along its sides we get a tessellation of $\matH^2$). 

We denote by $x,y,z$ both the sides of $\Delta$ and the reflections along them. For every $g\in \Gamma$ we define an abstract copy $\Delta_g$ of $\Delta$, and then we glue all these abstract copies altogether by pairing their sides  as follows: for every side of $\Delta$, say $x$, and every $g\in \Gamma$, we identify the two copies of the side $x$ in $\Delta_g$ and $\Delta_{gx}$ with the obvious identity map.

Via these identifications we get an abstract space $T$ tessellated into the triangles $\Delta_g$. The group $\Gamma$ acts freely and transitively on the tessellation (the element $g'\in \Gamma$ sends $\Delta_g$ to $\Delta_{g'g}$). We now prove that $T$ has a natural structure of a hyperbolic surface: since $\Gamma$ acts transitively, it suffices to check this for the points $p$ lying in $\Delta_e$. If $p$ lies in the interior of $\Delta_e$ or of a side, this is clear by construction. If $p$ is a vertex, say with inner angle $a$, by construction a cycle of $2a$ triangles $\Delta_e, \Delta_y, \Delta_{yz}, \Delta_{yzy}, \Delta_{yzyz}, \ldots $ is attached around $p$ because $(yz)^a = e$, and hence $T$ is naturally a hyperbolic surface also near $p$.

The hyperbolic surface $T$ is easily seen to be connected and complete. There is a natural developing map $\varphi\colon T\to \matH^2$ that sends $\Delta_g$ to $g(\Delta)$. The map $\varphi$ is a local isometry, hence a covering by Proposition \ref{local:isometry:prop}, hence an isometry since $\matH^2$ is simply connected. We identify $T$ with $\matH^2$ through $\varphi$.

It remains to prove that $\Gamma$ may be presented as stated. Let $G$ be the group presented as
$$\langle x, y,z\ |\ x^2, y^2, z^2, (xy)^c, (yz)^a, (zx)^b \rangle.$$
There is a natural surjection $G\to \Gamma$. We may repeat all the arguments above using $G$ instead of $\Gamma$ and get another $T'$ isometric to $\matH^2$ that covers $T$. We deduce that $T'=T$ and $G=\Gamma$.
\end{proof}

The triangle group $\Gamma(2,2,c) \isom \matZ/_{2\matZ}\times D_{2c}$ has order $4c$ and is the symmetry group of a prism. The triangle groups $\Gamma(2,3,3)$, $\Gamma(2,3,4)$, and $\Gamma(2,3,5)$ have order $24$, $48$, and $120$ and are the symmetry groups of the regular tetrahedron, octahedron, and icosahedron: the reader is invited to check all these facts visually by looking at Figure \ref{triangular:fig}.

The triangle groups $\Gamma(2,3,6)$, $\Gamma(2,4,4)$, and $\Gamma(3,3,3)$ are discrete groups of isometries of $\matR^2$ with compact quotient $\matR^2/_\Gamma$. 
If $\frac 1a + \frac 1b + \frac 1c <1$ the group $\Gamma(a,b,c)$ is a discrete subgroup of $\Iso(\matH^2)$. It contains infinitely many elliptic elements, such as reflections along lines and finite-order rotations around the vertices of the triangles. By Selberg's Lemma, there is a torsion-free subgroup $\Gamma' < \Gamma$ of some finite index $h$. The quotient $\matH^2/_{\Gamma'}$ is a closed hyperbolic surface and is tessellated into $h$ triangles isometric to $\Delta$.

\subsection{Coxeter polyhedra}
Triangle groups may be generalised to all dimensions as follows. A polyhedron $P$ in $\matH^n$ (or $\matR^n$, $S^n$) is a \emph{Coxeter polyhedron} if the dihedral angle of every codimension-two face divides $\pi$.\index{polyhedron!Coxeter polyhedron} For instance, the regular ideal tetrahedron and octahedron are Coxeter polyhedra.

The following theorem generalises Proposition \ref{triangle:gamma:prop}. Let $P$ be a finite Coxeter polyhedron: it is the convex hull of finitely many vertices in $\overline{\matH^n}$ (or $\matR^n$, $S^n$) and has $k$ facets, that we number as $1, \ldots, k$; we denote by $r_i$ the reflection along the $i$-th facet, and by $\frac {\pi}{a_{ij}}$ the dihedral angle formed by the $i$-th and $j$-th facets, if they are incident. Let $\Gamma$ be the isometry group generated by the reflections along the facets of $P$.

\begin{teo}
By mirroring iteratively a finite Coxeter polyhedron $P$ along its facets we get a tessellation of $\matH^n$ (or $\matR^n, S^n$).
The group $\Gamma$ acts freely and transitively on the tessellation: hence it is discrete and $P$ is a fundamental domain for $\Gamma$. A presentation for $\Gamma$ is
$$\langle r_1, \ldots, r_k \ |\  r_i^2, (r_ir_j)^{a_{ij}} \rangle$$
where $i$ varies in $1,\ldots, k$ and the pair $i,j$ varies among the incident facets.
\end{teo}
\begin{proof}
Same proof as Proposition \ref{triangle:gamma:prop}, with in addition an induction on the dimension $n$. Here are the details. We consider for simplicity only the hyperbolic case.

For every $g\in \Gamma$ we define an abstract copy $P_g$ of $P$, and we identify the $i$-th facet of $P_g$ and $P_{gr_i}$ for all $g\in \Gamma$ and all $i$. To prove that the resulting space $T$ is naturally a hyperbolic manifold, we use the induction hypothesis as follows. Every $p\in P_e$ lies in the interior of some $h$-dimensional face $f$, and let $\Gamma_f<\Gamma$ be the subgroup generated by the reflections along all the facets of $P$ containing $f$. The point $p$ is adjacent in $T$ to the polyhedra $P_g$ such that $g\in \Gamma_f$. By intersecting each such $P_g$ with a small codimension $h+1$ sphere centred in $p$ and contained in the codimension $h$ subspace orthogonal to $f$, we get some spherical polytope $Q_g$ of dimension $n-h-1<n$. (The spherical polytope $Q_t$ is usually called the \emph{spherical link} of $f$.)

By construction $Q_e \subset S^{n-h-1}$ is a Coxeter polytope and the subgroup $\Gamma_f < \Iso(S^{n-h-1})$ is generated by the reflections along its facets. By applying the induction hypothesis on $Q_e$ we deduce that the polytopes $Q_g$ with $g\in \Gamma_f$ form a tessellation of $S^{n-h-1}$, hence the polyhedra $P_g$ incident to $p$ form naturally a hyperbolic ball locally near $p$. Therefore $T$ is a hyperbolic manifold.

Proving that $T$ is complete requires a bit of care if $P_e$ has some ideal vertex $v$. In that case $v$ has a Euclidean link (obtained by intersecting $P_e$ with a small horosphere) and we conclude by induction as above that the Euclidean links of the $P_g$ incident to $v$ in $T$ glue to form a Euclidean $\matR^{n-1}$, that is a horosphere, so by intersecting each $P_g$ with small horoballs we get a horoball near $v$, which is complete.

The hyperbolic space $T$ is complete and connected, and we conclude as in the proof of Proposition \ref{triangle:gamma:prop}. 
\end{proof}

A group generated by some reflections along hyperplanes is called a \emph{reflection group}.\index{group!reflection group}
The following proposition shows that Coxeter polyhedra generate all the interesting reflection groups.

\begin{prop}
Every discrete reflection group $\Gamma$ is generated by the reflections along the facets of some Coxeter polyhedron.
\end{prop}
\begin{proof}
Consider the mirror hyperplanes of all the reflections in $\Gamma$. Since $\Gamma$ is discrete, these form a locally finite set and hence define a tessellation of $\matH^n$ onto which $\Gamma$ acts transitively. Pick one polyhedron $P$ of the tessellation. The reflections along the facets of $P$ generate $\Gamma$ (exercise).
\end{proof}

Coxeter polyhedra are beautiful objects that can be used to construct hyperbolic manifolds: every finite Coxeter polyhedron $P$ generates a reflection group $\Gamma$ which contains, by Selberg's Lemma, a torsion-free subgroup $\Gamma'$ of some finite index $h$. The quotient $M=\matH^n/_{\Gamma'}$ is a hyperbolic manifold and is tessellated into $h$ copies of $P$, so that $\Vol(M) = h\Vol(P)$. By residual finiteness, there are plenty of such manifolds.

\subsection{Coxeter graphs}
If a Coxeter polyhedron $P$ has a reasonable number of facets (for instance, if $P$ is a simplex), then one can study many of its properties by looking at its \emph{Coxeter graph},\index{Coxeter graph} which is constructed as follows: draw one node for each facet of $P$, and for every pair of facets intersecting with dihedral angle $\frac \pi a$ connect the corresponding nodes with an edge labeled with $a$. One also uses thickened or dashed edges to denote non-incident faces that do or do not intersect at some ideal vertex, respectively (these cases cannot occur on simplexes). Given the abundance of right angles, one usually omit the edges with label 2. 

\begin{ex}
Let $P$ be a regular polytope or tessellation with Schl\"afli symbol $\{a,b, \ldots, z\}$. By quotienting $P$ via all its isometries we get a Coxeter simplex, called its \emph{characteristic simplex}, whose Coxeter graph is\index{simplex!characteristic simplex}
\begin{center}
\includegraphics[width = 3 cm]{Coxeter_line}
\end{center}
\end{ex}

\subsection{Flat tori} \label{piatte:subsection}
The isometry group $\Iso(\matR^n)$ of $\matR^n$ contains the translations subgroup that we identify with $\matR^n$. The integer translations $\Gamma = \matZ^n$ form a discrete subgroup that acts freely on $\matR^n$, therefore the quotient $\matR^n /_\Gamma$ is a flat manifold. It is naturally diffeomorphic to the $n$-dimensional torus: 
$$\matR^n/_{\matZ^n} = \left(\matR/_{\matZ}\right)^n \isom \underbrace{S^1  \times \ldots \times S^1}_n.$$

\begin{ex} For every $p\in \matR^n$ the Dirichlet domain $D(p)$ is a $n$-dimensional unit cube centred at $p$.
\end{ex}

The flat $n$-torus may be seen as the unit $n$-cube with its opposite facets identified by translations. The two-dimensional case is shown in Figure \ref{domini:fig}-(left): by identifying the opposite sides of a square we get a torus. \index{flat torus}

\begin{figure}
\begin{center}
\includegraphics[width = 11.5 cm] {\iftoggle{BW}{domini-BW}{domini}}
\nota{A fundamental domain in $\matR^2$ for the torus (left) and the Klein bottle (centre): opposite sides should be identified as indicated by the arrows. A fundamental domain in $S^2$ for $\matRP^2$ (right).}
\label{domini:fig}
\end{center}
\end{figure}

The $n$-torus possesses a continuous family of non-isometric flat metrics. A \emph{lattice} $\Gamma<\matR^n$ is a discrete subgroup isomorphic to $\matZ^n$ which spans $\matR^n$ as a vector space. We see $\Gamma$ as a group of translations. \index{lattice}

\begin{ex} For every lattice $\Gamma$, the flat manifold $\matR^n/_\Gamma$ is diffeomorphic to the $n$-torus. A fundamental domain is the parallelotope spanned by $n$ generators of $\Gamma$.
\end{ex}

A Dirichlet domain for $\Gamma$ is almost never a parallelotope! In the \emph{hexagonal torus} $\matR^2/_\Gamma$ the group $\Gamma$ is the equilateral lattice generated by the translations\index{torus!hexagonal torus}
$$(1,0), \quad \big(\tfrac 12, \tfrac{\sqrt 3}2\big)$$ 
and the Dirichlet domain of any point is a regular hexagon. Orientable flat manifolds are easily classified (up to diffeomorphism) in dimension two. \index{hexagonal torus}

\begin{prop} \label{flat:is:torus:prop}
Every closed flat orientable surface is a torus.
\end{prop}
\begin{proof}
A closed flat surface $S$ is isometric to $\matR^2/_\Gamma$ for some discrete group $\Gamma$ of orientation-preserving isometries acting freely. Every fixed-point-free orientation-preserving isometry of $\matR^2$ is a translation (exercise), hence $\Gamma < \matR^2$. If $\Gamma$ has rank one then $S$ is not compact, hence $\Gamma$ is a lattice and $S$ is a torus.
\end{proof}

Flat manifolds more complicated than tori can be constructed by considering also non-translational isometries of $\matR^n$, which exist for $n\geqslant 3$ and also for $n=2$ in the non-orientable setting. For instance, in dimension two we find the \emph{Klein bottle} by taking $\Gamma$ as the group generated by the isometries:
$$\tau\colon(x,y) \mapsto (x+1, y), \qquad \eta\colon(x,y) \mapsto (1-x, y+1).$$
A fundamental domain for the Klein bottle is shown in Figure \ref{domini:fig}-(centre).
\index{Klein bottle}

Every subgroup $\Gamma' < \Gamma$ furnishes another flat manifold $\matR^2/_{\Gamma'}$ that covers the Klein bottle. For instance, the subgroups $\langle \tau \rangle$ and $\langle \eta \rangle$ generated respectively by $\tau$ and $\eta$ are both isomorphic to $\matZ$ but provide two different coverings: the manifold $\matR^2/_{\langle\tau\rangle}$ is an infinite cylinder while $\matR^2/_{\langle\eta\rangle}$ is an infinite M\"obius strip. 

The subgroup $\Gamma'$ generated by the translations $\tau$ and $\eta^2$ is isomorphic to $\matZ^2$ and has index 2 in $\Gamma$. The Klein bottle is doubly covered by the flat torus $\matH^2/_{\Gamma'}$. A fundamental domain for $\Gamma'$ is a rectangle with vertices $(0,0)$, $(1,0)$, $(0,2)$, $(1,2)$.

We have seen that the Klein bottle is covered by a torus. In fact, we will see in Section \ref{Bieberbach:section} that \emph{every} closed flat $n$-manifold is covered by a flat $n$-torus.

\begin{ex}
A square $Q\subset \matR^2$ is a Coxeter polygon and determines a reflection group $\Gamma$. What is the minimum index of a torsion-free subgroup $\Gamma' < \Gamma$? Which topological surfaces $\matR^2/_{\Gamma'}$ do we get?
\end{ex}

\subsection{Real projective spaces} 
We now exhibit some elliptic manifolds. Every elliptic manifold is covered by $S^n$ and is hence compact and has finite fundamental group (because coverings between compact manifolds have finite degree). \index{$\matRP^n$}

An important example is the real projective space $\matRP^n = S^n/_\Gamma$ where $\Gamma$ is the order-two group generated by the antipodal map $\iota (x) = -x$. 

\begin{ex}
For every $p\in S^n$, the Dirichlet domain $D(p)$ of $\Gamma$ is the hemisphere centred at $p$.
\end{ex}
The two-dimensional case is shown in Figure \ref{domini:fig}-(right). The only elliptic surfaces are $S^2$ and $\matRP^2$ in virtue of the following.
\begin{prop}
In even dimension $n$ the only elliptic manifolds are $S^n$ and $\matRP^n$.
\end{prop}
\begin{proof}
Let $M=S^n/_\Gamma$ be an elliptic manifold. Every matrix in $\SO(n+1)$ has an eigenvalue $+1$ because $n$ is even (exercise) and hence acts in $S^n$ with a fixed point. Therefore $\Gamma$ contains no non-trivial orientation-preserving isometries, and since these form a group of index at most two, either $\Gamma$ is trivial or is generated by an orientation-reversing fixed-point free involution in $O(n+1)$, and $-I$ is the only such element (exercise).
\end{proof}

\subsection{Lens spaces} \label{lens:subsection}
In dimension two $\matRP^2$ and $S^2$ are the only elliptic manifolds, but in dimension three there are more. Let $p\geqslant 1$ and $q\geqslant 1$ be coprime integers and set $\omega = e^{\frac{2\pi i}p}$. We identify $\matR^4$ with $\matC^2$ and see $S^3$ as
$$S^3 = \big\{(z,w)\in \matC^2\ \big|\ |z|^2+|w|^2 = 1 \big\}.$$
The map
$$f(z,w) = (\omega z, \omega^qw)$$
is an isometry of $\matR^4$ because it consists of two simultaneous rotations on the coordinate planes $w=0$ and $z=0$. The map $f$ hence induces an isometry of $S^3$. It has order $p$ and none of its iterates $f, f^2, \ldots, f^{p-1}$ has a fixed point. Therefore the group $\Gamma = \langle f \rangle$ generated by $f$ acts freely on $S^3$, and is discrete because it is finite. \index{lens space} \index{$L(p,q)$}

We have constructed an elliptic manifold $S^3/_\Gamma$, called a \emph{lens space} and indicated with the symbol $L(p,q)$. Its fundamental group is isomorphic to $\Gamma = \matZ/_{p\matZ}$. Note that the manifold depends on both $p$ and $q$.

\section{Geodesic boundary, non-complete, and cone manifolds}
We now introduce three important variations on the complete hyperbolic manifolds theme.
The first consists of admitting a totally geodesic boundary, the second is a brief overlook of some phenomena that may occur when the completeness hypothesis is dropped, and in the third we allow some conical singularity on a codimension-two geodesic stratum.

\subsection{Hyperbolic manifolds with geodesic boundary}
We reformulate a definition of hyperbolic (elliptic, flat) manifolds that allows the presence of some geodesic boundary. These manifolds are useful because they can be glued along their boundaries to produce new hyperbolic (elliptic, flat) manifolds.
\index{hyperbolic manifold!with geodesic boundary}
\index{flat manifold!with geodesic boundary}
\index{elliptic manifold!with geodesic boundary}
\begin{defn} A \emph{hyperbolic} (\emph{elliptic, flat}) \emph{manifold} $M$ \emph{with geodesic boundary}
is a Riemannian manifold with boundary where every point has an open neighbourhood isometric to an open set in a half-space in $\matH^n$ ($S^n$, $\matR^n$).
\end{defn}

The boundary $\partial M$ of a hyperbolic (elliptic, flat) $n$-manifold with geodesic boundary is a hyperbolic (elliptic, flat) $(n-1)$-manifold without boundary. Theorem \ref{simply:teo} extends appropriately to this context.

\begin{teo}
Every complete simply connected hyperbolic $n$-manifold $M$ with geodesic boundary is isometric to the intersection of some half-spaces in $\matH^n$ with disjoint boundaries.
\end{teo}
\begin{proof}
We construct a developing map $D\colon M\to\matH^n$ as in Theorem \ref{simply:teo}. The map $D$ is injective, because every two points $p,q \in M$ are connected by a geodesic (exercise: use completeness and geodesic boundary), which is sent to a geodesic via $D$, so $p,q$ are sent to distinct points $D(p), D(q)$. 

Therefore $D(M)\subset \matH^n$ is isometric to $M$. Since it is complete and with geodesic boundary, its boundary consists of disjoint hyperplanes, and hence $D(M)$ is the intersection of half-spaces bounded by them.  
\end{proof}

An example is sketched in Figure \ref{universale:fig}. Every complete hyperbolic manifold with geodesic boundary can be enlarged to a hyperbolic manifold without boundary in a canonical way.

\begin{figure}
\begin{center}
\includegraphics[width = 4cm] {\iftoggle{BW}{universale-BW}{universale}}
\nota{An intersection of (possibly infinitely many!) half-planes. The universal cover of a hyperbolic surface with boundary is isometric to such an object.}
\label{universale:fig}
\end{center}
\end{figure}


\begin{cor} Every complete hyperbolic $n$-manifold $M$ with geodesic boundary is contained in a complete hyperbolic $n$-manifold $N$ without boundary, such that $N\setminus \interior M$ is diffeomorphic to $\partial M \times [0,+\infty)$.
\end{cor}
\begin{proof}
The proof of Proposition \ref{corto:prop} applies and shows that $M = \tilde M/_\Gamma$ where $\tilde M\subset \tilde \matH^n$ is an intersection of half-spaces and $\Gamma< \Iso(\tilde M)$ acts freely and properly discontinuously on $\tilde M$.

Every local isometry in $\matH^n$ extends to a global isometry, therefore $\Gamma < \Iso(\tilde M) < \Iso(\matH^n)$. The group $\Gamma$ acts freely on $\tilde M$ and hence also on $\matH^n$: if it had a fixed point $x\in \matH^n$, it would fix also the unique point $y \in \tilde M$ that is the closest to $x$.
The manifold $N=\matH^n/_\Gamma$ contains naturally $M$.

For every $p \in \partial M$, let $\gamma_p(t)$ be the unit speed geodesic starting at $p$ orthogonal to $\partial M$ and directed outward. By looking at the universal cover one sees easily that the map $\partial M \times [0,+\infty) \to N \setminus \interior M$, $(p,t)\mapsto \gamma_p(t)$ is a diffeomorphism (we use here that $\tilde M\subset \matH^n$ is convex).
\end{proof}

\subsection{Cut and paste}
Hyperbolic manifolds with geodesic boundary are useful because they can be glued to produce new hyperbolic manifolds.

Let $M$ and $N$ be hyperbolic manifolds with geodesic boundary and $\psi\colon \partial M \to \partial N$ be an isometry. Let $M\cup_\psi N$ be the topological space obtained by quotienting the disjoint union $M\sqcup N$ by the equivalence relation that identifies $p$ to $\psi(p)$ for all $p\in \partial M$.

\begin{prop} \label{geodesic:glue:prop}
The space $M\cup_\psi N$ has a natural structure of hyperbolic manifold.
\end{prop}
\begin{proof}
The interiors of $M$ and $N$ inherit their hyperbolic metrics. When we glue $p$ to $\psi(p)$ we get a point $q$ that has two half-disc neighbourhoods on its sides, which glue to a honest hyperbolic disc, inducing a hyperbolic metric near $q$.
\end{proof}

Conversely, if an orientable hyperbolic manifold $M$ contains an orientable geodesic hypersurface $N$, we can cut $M$ along $N$ to get a hyperbolic manifold with geodesic boundary. The boundary will consist of two copies of $N$.

\subsection{Non-complete hyperbolic manifolds} \label{non-complete:subsection}
There is no classification of simply connected non-complete hyperbolic manifolds: for instance, we may get plenty of uninteresting examples by removing complicated closed sets from $\matH^n$. However, the first part of the proof of Theorem \ref{simply:teo} still applies and provides the following:

\begin{prop}
Let $M$ be a non-complete simply connected hyperbolic $n$-manifold. There is a local isometry
$$D\colon M \to \matH^n$$
which is unique up to post-composing with $\Iso(\matH^n)$.
\end{prop}
\begin{proof}
Construct $D$ as in the proof of Theorem \ref{simply:teo}: the completeness of $M$ is used there only in the last paragraph to show that $D$ is a covering.
As a local isometry, the map $D$ is determined by its first-order behaviour at any point $p\in M$, and is hence unique up to post-composing with an isometry of $\matH^n$.
\end{proof}

The map $D$ is called a \emph{developing map} and is neither injective nor surjective in general. We can push-forward isometries along $D$ as follows. \index{developing map}

\begin{prop}
Let $M$ be a hyperbolic $n$-manifold and $D\colon M \to \matH^n$ be a local isometry. For every $g\in \Iso(M)$ there is a unique $\rho(g)\in \Iso(\matH^n)$ such that
$$\rho(g) \circ D = D\circ g.$$
The resulting map $\rho \colon \Iso(M) \to \Iso(\matH^n)$ is a homomorphism.
\end{prop}
\begin{proof}
Pick a point $p\in M$ and define $\rho(g)$ as the unique isometry of $\matH^n$ such that $\rho(g) \circ D = D\circ g$ on a small neighbourhood of $p$. This equality between local isometries holds locally and hence globally on $M$. The map $\rho$ is easily checked to be a homomorphism.
\end{proof}

The homomorphism $\rho$ is the \emph{holonomy} associated to $D$. If $M$ is a non-complete hyperbolic $n$-manifold, its universal covering $\pi\colon \tilde M \to M$ inherits a non-complete hyperbolic metric. Therefore we get a developing map 
$$D\colon \tilde M \to \matH^n$$ 
together with a holonomy \index{holonomy}
$$\rho \colon \Aut(\pi) \longrightarrow \Iso(\matH^n)$$
which is the restriction of $\rho$ to the subgroup $\Aut(\pi)<\Iso(\tilde M)$. By fixing a point in $\tilde M$ we identify $\Aut(\pi)$ with $\pi_1(M)$ and get a holonomy
$$\rho\colon \pi_1(M) \longrightarrow \Iso(\matH^n).$$

As every metric space, a non-complete Riemannian manifold $M$ has a unique \emph{completion} $\overline M$, a complete metric space that contains $M$ as an open dense set. The completion $\overline M$ is however not necessarily a manifold, except in some lucky cases. These lucky cases are extremely important in dimension three, as we will see in Chapter \ref{three:chapter}. In that chapter we will understand the completion $\overline M$ of a hyperbolic three-manifold by studying its developing map and holonomy.

\subsection{Singularities} \label{singularities:subsection}
A particular class of non-complete hyperbolic manifolds deserves our attention, because their completions are some nice and natural objects: manifolds with \emph{cone angles} along some codimension-two singular geodesic stratum.

Let $S\subset \matH^n$ be a codimension-two subspace. The incomplete hyperbolic manifold $\matH^n \setminus S$ has fundamental group $\matZ$ and we denote by
$$\pi \colon X \longrightarrow \matH^n \setminus S$$
its universal covering: here $X$ is an interesting non-complete simply connected hyperbolic manifold. Note that $\pi$ may be interpreted as a developing map $\pi\colon X\to \matH^n$ which is neither injective nor surjective.

\begin{ex}
The metric completion of $X$ is $\overline X = X \sqcup \tilde S$ where $\tilde S$ is an identical copy of $S$. The covering $\pi$ extends to a surjective map
$$\pi\colon \overline X \to \matH^n$$
that sends $\tilde S$ to $S$. If $p\in \tilde S$ and $q\in \overline X$ then $d(p,q) = d(\pi(p), \pi(q)).$
\end{ex}

The space $\tilde S$ is the \emph{singularity} of $\overline X$ and should be interpreted as lying in $\overline X$ with infinite cone angle: the space $\overline X$ looks like a book with infinitely many pages (it is not locally compact at $\tilde S$). Singularities with finite cone angles will be defined in the next section as quotients of this model. The dimensions that will be interesting for us are of course $n=2$ and $3$, where $\tilde S$ is a point and a line, respectively.

\subsection{Cone manifolds}
For every $\theta\in \matR$ there is a well-defined rotation $\Rot_\theta\colon \overline X \to \overline X$ of angle $\theta$ around $\tilde S$, which projects to a rotation in $\matH^n$ of angle $\theta$ around $S$; note that $\Rot_\theta \neq \id$ for all $\theta \neq 0$, including $\theta = 2k\pi$.

We fix $\theta \neq 0$, so that $\Rot_\theta$ generates a free cyclic group $\Gamma < \Iso(\overline X)$.

\begin{ex}
The metric space $\overline X/_\Gamma$ is the completion of $X/_\Gamma$ and is homeomorphic to $\matR^n$.
\end{ex}

The image of $\tilde S$ along the quotient $\pi\colon \overline X \to \overline X /_\Gamma$ is a copy of $\tilde S$ and should now be interpreted as a singularity with cone angle $\theta$. When $\theta = 2k\pi$ for some integer $k$ the map $\pi\colon \overline X \to \matH^n$ factors as
$$\overline X \overset {\pi_1} {\longrightarrow} \overline X/_\Gamma \overset {\pi_2} \longrightarrow \matH^n.$$
When $k=\pm 1$ the map $\pi_2$ is an isometry. We introduce the following definition.\index{hyperbolic manifold!with cone angles}

\begin{defn}
A \emph{hyperbolic manifold $M$ with cone angles} is a metric space with charts in some $\overline X/_\Gamma$ and transition maps that are isometries.
\end{defn}

The term ``isometries'' should be interpreted in the strongest sense: the transition maps send singular points to singular points, and outside of the singular points these are isometries of Riemannian manifolds.

The points in $M$ that are mapped along charts to some singular set form a hyperbolic codimension-two submanifold in $M$ with some cone angle $\theta$, and different components may have different cone angles.
Points with cone angle $2\pi$ may be considered as ordinary points, while points with cone angle different from $2\pi$ are \emph{singular} and form the \emph{singular locus} of $M$.

The singular locus of $M$ is a geodesic codimension-two manifold, and its complement is an ordinary hyperbolic manifold, whose metric completion is $M$.
Flat and spherical manifolds with cone angles are defined in the same way. \index{elliptic manifold!with cone angles} \index{flat manifold!with cone angles}

\subsection{Examples}
The following simple construction is already a source of many two-dimensional examples: let $P$ be any polygon in $\matH^2, \matR^2,$ or $S^2$, with some inner angles $\alpha_1,\ldots, \alpha_k$. By doubling $P$ along its edges we get a topological sphere with $k$ cone points of angles $2\alpha_1, \ldots, 2\alpha_k < 2\pi$. This is a hyperbolic, flat, or spherical surface with cone angles.

Another simple but maybe more intriguing construction is the following: pick any polyhedron $P$ in $\matH^3, \matR^3$, or $\matS^3$. The faces of $P$ are (hyperbolic, flat, or elliptic) polygons that glue isometrically along the edges: hence $\partial P$ is a topological sphere with a natural structure of (hyperbolic, flat, or elliptic) manifold with cone angles at the vertices.

To construct examples in dimension three we need a bit more work because polyhedra have more strata. Non-compact examples are easier to build: let $P\subset \matH^3$ be any ideal hyperbolic polyhedron, with some $k$ edges with dihedral angles $\alpha_1,\ldots, \alpha_k$; the double of $P$ is naturally a hyperbolic cone manifold where the edges form the singular locus with cone angles $2\alpha_1, \ldots, 2\alpha_k$.

\section{Orbifolds} 
An orbifold is an object locally modelled on finite quotients of $\matR^n$. It naturally arises when we quotient a Riemannian manifold by a discrete group of isometries that may not act freely. Orbifolds behave like manifolds on many aspects.

\subsection{Definition} Let $\Gamma < \On(n)$ be a finite group of linear isometries and $V\subset \matR^n$ be a $\Gamma$-invariant open set. The resulting map
$$\varphi \colon V \longrightarrow V/_\Gamma$$
is called a \emph{local orbifold model}. For instance, we may pick $V=B(0,r)$. \index{orbifold!local model of an orbifold}

We now generalise manifolds by allowing local orbifold models in the atlas.\index{orbifold}
\begin{defn} Let $O$ be a Hausdorff paracompact topological space. An \emph{orbifold atlas} on $O$ is an open covering $\{U_i\}_{i\in I}$ of $O$, closed by finite intersections and equipped with local orbifold models 
$$\varphi_i\colon V_i \longrightarrow V_i/_{\Gamma_i} = U_i.$$
The local models are connected by some appropriate transition functions: for every inclusion $U_i\subset U_j$ there is an injective homomorphism 
$$f_{ij}\colon \Gamma_i \hookrightarrow \Gamma_j$$ 
and a $\Gamma_i$-equivariant smooth embedding $\psi_{ij} \colon V_i \hookrightarrow V_j$ compatible with the local models, that is 
$$\varphi_j \circ \psi_{ij} = \varphi_i.$$
\end{defn}

Two such atlases are equivalent if the are contained in some bigger atlas. The space $O$ equipped with an orbifold atlas is an \emph{orbifold} of dimension $n$.
See an example in Figure \ref{orbifold_charts:fig}. 

\begin{figure}
\begin{center}
\includegraphics[width = 8 cm] {\iftoggle{BW}{orbifold_charts-BW}{orbifold_charts}}
\nota{An orbifold $O$ with three local models $O=U_1 \supset U_2 \supset U_3$. The group $\Gamma_1 = \matZ_2\times \matZ_2$ acts on $V_1$ by reflections on the two coordinate axis, the group $\Gamma_2 = \matZ_2$ acts on $V_2$ by reflection on the horizontal axis, the group $\Gamma_3$ is trivial.}
\label{orbifold_charts:fig}
\end{center}
\end{figure}

\begin{oss}
We think at the maps $\psi_{ij}$ and $f_{ij}$ as defined only up to the action of $\Gamma_j$ (which acts on $\psi_{ij}$ by composition and on $f_{ij}$ by conjugation).  In particular, if $U_i\subset U_j \subset U_k$ then we can verify that the equalities $\psi_{ik} = \psi_{jk}\circ \psi_{ij}$ and $f_{ik} = f_{jk}\circ f_{ij}$ hold only up to this ambiguity. See Figure \ref{orbifold_charts:fig}.
\end{oss}
The \emph{isotropy group} $\Gamma_p$ of a point $p\in O$ is the stabiliser of any lift of $p$ in any local model $V$ with respect to the action of $\Gamma$. By definition $\Gamma_p$ is a finite subgroup of $\On(n)$. A point $p$ is \emph{regular} if its isotropy group is trivial, and \emph{singular} otherwise. \index{orbifold!isotropy group of an orbifold}

\begin{example} 
The isotropy group $O_p$ in the orbifold $O$ from Figure \ref{orbifold_charts:fig} is 
$\matZ_2 \times \matZ_2$ at the origin, $\matZ_2$ at the axis, and trivial elsewhere. 
\end{example}

\begin{prop}
The regular points in an orbifold form a dense subset. 
\end{prop}
\begin{proof}
On a local model $V \to V/_\Gamma$, a singular point is the image of the fixed point locus of some element in $\Gamma$, which is in turn a proper subspace.
\end{proof}

An orbifold is \emph{locally oriented} if all the finite groups $\Gamma$ lie in $\SO(n)$, and it is \emph{oriented} if, in addition, all the transition maps $\psi_{ij}$ are orientation-preserving. In a locally oriented orbifold reflections are not admitted and hence the singular locus has codimension $\geqslant 2$.
An open subset of an orbifold is naturally an orbifold.

\subsection{Examples} Orbifolds are natural objects and there are plenty of nice examples around.

\begin{example} A differentiable manifold is an orbifold whose points are all regular. A differentiable manifold with boundary may be interpreted as an orbifold whose boundary points have the local structure of type $\matR^n/_\Gamma$ where $\Gamma = \matZ_2$ is generated by a reflection along a hyperplane. The boundary should be interpreted as a \emph{mirror}.
\end{example}

\begin{ex}
Construct an orbifold structure on the triangle as suggested by Figure \ref{orbifold_triangle:fig}.
\end{ex}

\begin{example}
Let $\Gamma < \On(n)$ be finite and $V\subset \matR^n$ a $\Gamma$-invariant open set. The quotient $V/_\Gamma$ has an orbifold structure, defined by the unique local model $V \to V/_\Gamma$. 
\end{example}

\begin{figure}
\begin{center}
\includegraphics[width = 2.2 cm] {\iftoggle{BW}{orbifold_triangle-BW}{orbifold_triangle}}
\nota{An orbifold structure on the triangle constructed using the local models from Figure \ref{orbifold_charts:fig}. The isotropy group is $\matZ_2 \times \matZ_2$ at the vertices, $\matZ_2$ at the edges, and trivial in the interior.}
\label{orbifold_triangle:fig}
\end{center}
\end{figure}

We generalise the last example.

\begin{prop}
If $M$ a Riemannian manifold and $\Gamma<\Iso(M)$ is a discrete subgroup, the quotient $M/_\Gamma$ has a natural orbifold structure.
\end{prop}
\begin{proof}
Take a point $p\in M/_\Gamma$ and $\tilde p\in M$ a lift. Since $\Gamma$ is discrete, the stabiliser $\Gamma_{\tilde p}$ of $\tilde p$ is finite and there is a $r>0$ such that $\exp_{\tilde p}(B_r(0)) = B_r(\tilde p)$ and $g(B_r(\tilde p))$ intersects $B_r(\tilde p)$ for some $g\in \Gamma$ if and only if $g\in\Gamma_{\tilde p}$. 

The ball $B_r(\tilde p)$ is clearly $\Gamma_{\tilde p}$-invariant. The group $\Gamma_{\tilde p}$ acts linearly and orthogonally on $B_r(0)$ and we get an orbifold local model
$$B_r(0) \longrightarrow B_r(0) /_{\Gamma_{\tilde p}} \isom B_r(\tilde p)/_{\Gamma_{\tilde p}} = U_p$$
where the diffeomorphism is induced by the exponential map. We have constructed an orbifold local model $U_p$ around $p$. We extend the covering $\{U_p\}$ thus obtained by adding all the non-empty intersections. 
\end{proof}

This is the richest source of nice examples. The quotient map $M \to M/_\Gamma$ is a covering, in an appropriate sense that we now explain.

\subsection{Coverings}
Let $\Gamma' < \Gamma < \On(n)$ be finite groups and $V\subset \matR^n$ a $\Gamma$-invariant open set. The natural map
$$\varphi\colon V/_{\Gamma'} \longmapsto V/_{\Gamma}$$
between the two orbifolds is a \emph{local covering}. \index{orbifold!orbifold covering}

\begin{defn}
A continuous map $p\colon \tilde O \to O$ between orbifolds is a \emph{covering} if every point $p\in O$ has a neighbourhood $U$ with $p^{-1}(U) = \sqcup_{i \in I} U_i$ and every restriction $p|_{U_i}\colon U_i \to U$ is a local covering.
\end{defn}

The following is the main source of examples.

\begin{example}
Let $M$ be a Riemannian manifold and $\Gamma' < \Gamma < \Iso(M)$ be discrete groups. The natural map $M/_{\Gamma'} \to M/_\Gamma$ is an orbifold covering.
\end{example}

\begin{defn}[The good, the bad, and the very good]
An orbifold is \emph{good} if it is covered by a manifold, and it is \emph{bad} otherwise. It is \emph{very good} if it is finitely covered by a manifold. \index{orbifold!the bad, the good, and the very good orbifold}
\end{defn}

\subsection{Hyperbolic, flat, and elliptic orbifolds} \label{more:examples:subsection}
We define a \emph{hyperbolic, flat}, or \emph{elliptic} orbifold to be an orbifold whose local models are isometric quotients of balls in $\matH^n$, $\matR^n$, or $S^n$, and whose transition functions are also isometries. \index{orbifold!hyperbolic, flat, elliptic orbifold}

The quotient of a hyperbolic, flat, or spherical manifold by a discrete group of isometries is naturally a hyperbolic, flat, or elliptic orbifold.

\begin{example}
The triangle group $\Gamma(a,b,c)$ introduced in Section \ref{triangle:groups:subsection} is a discrete group of isometries of $S^2, \matR^2$, or $\matH^2$. It defines a triangular orbifold $\Delta$, which is hyperbolic, flat, or elliptic, according to whether $\frac 1a+\frac 1b + \frac 1c$ is smaller, equal, or bigger than 1. The isotropy groups are the dihedral $D_{2a}, D_{2b}$, and $D_{2c}$ at the vertices, $\matZ_2$ at the edges, and trivial elsewhere.

The index-two subgroup $\Gamma^{\rm or}(a,b,c) \triangleleft \Gamma(a,b,c)$ consisting of orientation-preserving isometries is sometimes called a \emph{von Dyck group}. It defines an orientable orbifold $O$ that double-covers $\Delta$. The orbifold $O$ is a topological sphere with three singular points with rotation isotopy groups $\matZ/_{a\matZ}, \matZ/_{b\matZ}, \matZ/_{c\matZ}$. See Figure \ref{orbifold_triangle_abc:fig}.
\index{group!von Dyck group}
\end{example}

\begin{ex} \label{VD:ex}
Show that the following is a presentation for $\Gamma^{\rm or}(a,b,c)$:
$$\langle\ r, s, t\ |\ r^{a}, s^{b}, t^{c}, rst\ \rangle.$$
The triple $r,s,t$ of generators is intrinsically determined in the group up to simultaneous conjugation or inversion. Different unordered triples $(a,b,c)$ produce non-isomorphic Von Dyck groups $\Gamma^{\rm or}(a,b,c)$.
\end{ex}
\begin{proof}[Hints]
The group $\Gamma^{\rm}(a,b,c)$ preserves the tessellation $T$ of $\matH^2$, $\matR^2$, or $S^2$ into triangles with angles $\frac \pi a, \frac \pi b, \frac \pi c$, so its finite-order elements are rotations along some vertices of $T$. Three rotations $r,s,t$ that satisfy $rst=1$ are of some very special kind.
\end{proof}

\begin{example}
More generally, every hyperbolic (or flat, elliptic) Coxeter polyhedron is a hyperbolic (or flat, elliptic) orbifold.
\end{example}

\begin{example}
Quotient the flat torus $T = \matR^2/_{\matZ^2}$ by the \emph{elliptic involution} $(x,y)\mapsto (-x,-y)$. This isometry of $T$ has four fixed points and the quotient flat orbifold is a sphere with four singular points, see Figure \ref{elliptic_involution:fig}.
\end{example}

\begin{figure}
\begin{center}
\includegraphics[width = 8 cm] {\iftoggle{BW}{orbifold_triangle_abc-BW}{orbifold_triangle_abc}}
\nota{The index-two orientation-preserving $\Gamma^{\rm or}(a,b,c) \triangleleft \Gamma(a,b,c)$ defines an index-two orbifold covering $O \to \Delta$, where $O$ is a 2-sphere with three singular points and $\Delta$ a triangle.}
\label{orbifold_triangle_abc:fig}
\end{center}
\end{figure}

\begin{figure}
\begin{center}
\includegraphics[width = 9 cm] {\iftoggle{BW}{elliptic_involution-BW}{elliptic_involution}}
\nota{The elliptic involution quotients the torus to a sphere with four singular points. The fixed points and their images are drawn in \iftoggle{BW}{grey}{red}. We show both a planar (left) and spacial (right) picture.}
\label{elliptic_involution:fig}
\end{center}
\end{figure}

\begin{figure}
\begin{center}
\includegraphics[width = 10 cm] {\iftoggle{BW}{orbifold_Farey-BW}{orbifold_Farey}}
\nota{The hyperbolic orbifold $\matH^2/_{\PSLZ}$ is obtained by mirroring the sides of a triangle with inner angles $\frac \pi 3$, $\frac \pi 2$, and zero (left). The orbifold is topologically a punctured sphere with two singular points with rotational isotropy $\matZ_2$ and $\matZ_3$ (right).}
\label{orbifold_Farey:fig}
\end{center}
\end{figure}

There are also many interesting non-compact orbifolds.

\begin{example}
A fundamental domain for $\PSLZ$ was shown in Figure \ref{modular_group:fig}. By gluing its sides according to the action we see quite easily that the hyperbolic orbifold $H^2/_{\PSLZ}$ is non-compact and contains two rotational points of order 2 and 3, see Figure \ref{orbifold_Farey:fig}.
\end{example}

\begin{ex} The 1-dimensional compact orbifolds are $S^1$ and the segment $[0,1]$ with mirrored endpoints. Show that there are coverings $S^1 \to [0,1]$ of degree 2, and also coverings $[0,1] \to [0,1]$ of any positive degree.
\end{ex}

When $O=X/_\Gamma$ and $X$ is a simply connected manifold, we say that $\Gamma$ is the \emph{fundamental group} of $O$. It is also possible to define the fundamental group for a generic orbifold using the appropriate orbi-notions of paths and homotopies. 

\subsection{Cone manifolds}
Hyperbolic (flat, elliptic) cone manifolds and orbifolds are different objects, but they have a wide common intersection.

A hyperbolic (flat, elliptic) orbifold $O$ whose isotropy groups $O_p$ are either trivial or generated by a $\frac{2\pi }p$-rotation along a codimension-two subspace is naturally a hyperbolic (flat, elliptic) cone manifold with cone angles $\frac{2 \pi}p$. The singular set may consist of various connected components and the natural number $p$ depends on the component. 

Conversely, every hyperbolic (flat, elliptic) manifold with cone angles that divide $2\pi$ can be given a natural hyperbolic (flat, elliptic) orbifold structure whose singular set consists of rotational points only. Moreover, we can easily prove that it is a good one:

\begin{prop} \label{from:cone:to:orbifold:prop}
Every hyperbolic (flat, elliptic) manifold $M$ with cone angles that divide $2\pi$ is a good hyperbolic (flat, elliptic) orbifold. If $M$ is complete, then $M=\matH^n/_\Gamma$ ($\matR^n/_\Gamma$, $S^n/_\Gamma$) for some discrete group $\Gamma$ of isometries.
\end{prop}
\begin{proof}
We consider for simplicity only the hyperbolic case. We must prove that $M$ is orbifold-covered by a hyperbolic manifold.

The singular set $S\subset M$ decomposes into connected components $S = \sqcup_i S_i$, each with a cone angle $\frac{2\pi}{p_i}$ for some integer $p_i \geqslant 2$. The complement $M' = M \setminus S$ is a non-complete hyperbolic manifold and as such it is equipped with a developing map $D\colon \tilde M' \to \matH^n$ and a holonomy $\rho\colon \pi_1(M') \to \Iso(\matH^n)$. 

We pick a loop $\mu_i$ around each component $S_i$ of $S$ and note that by definition $\rho(\mu_i)$ is an elliptic isometry that rotates $\matH^n$ around a codimension-two subspace by the angle $\frac{2\pi}{p_i}$. In particular $\rho(\mu_i)$ has order $p_i$. 

We consider the regular covering $N \to M'$ corresponding to the subgroup $\ker \rho$ and give $N$ the hyperbolic structure induced by $M'$. The completion $\overline N$ of $N$ is a hyperbolic manifold that orbifold-covers $M$: since $\rho(\mu_i)$ has order precisely $p_i$, the added singular points in $\overline N$ have cone angle $2\pi$ and are hence ordinary. 

If $M$ is complete, then $\overline N$ also is and hence $\overline N = \matH^n/_{\Gamma'}$. We define $\Gamma<\Iso(\matH^n)$ to be the group generated by $\Gamma'$ and any lifts of the deck transformations of $\overline N\to M$, and we get $M=\matH^n/_\Gamma$.
\end{proof}

\begin{cor} \label{very:good:cor}
If $M$ is complete and has finitely generated fundamental group, then it is very good. 
\end{cor}
\begin{proof}
We know that $M=\matH^n/_\Gamma$ and we apply Selberg's Lemma.
\end{proof}

Summing up, complete hyperbolic cone manifolds with angles that divide $2\pi$ can be promoted to good orbifolds and hence to quotients $\matH^n/_\Gamma$. We will encounter many examples in dimensions two and three later on.

\subsection{References}
All the material presented here is standard and introduced in various books, like Benedetti -- Petronio \cite{BP} and Ratcliffe \cite{R}. A beautiful introduction to regular polytopes is Coxeter's 1963 book \cite{Cox}, while the reader interested in Coxeter polytopes may consult Vinberg's survey \cite{Vin}.

The theory of orbifolds and cone manifolds presented here is quite limited: to get more background, the reader may consult Thurston's notes \cite[Chapter 13]{Th} for the orbifolds, and McMullen's paper \cite{McM-cone} for the cone manifolds.


